% ============================================================
%  Технический отчёт по проекту «Секретный Архив»
% ============================================================
\documentclass[12pt, a4paper]{article}

\usepackage[utf8]{inputenc}
\usepackage[T2A]{fontenc}
\usepackage[russian]{babel}
\usepackage[top=2.5cm, bottom=2.5cm, left=3cm, right=2cm]{geometry}
\usepackage{xcolor}
\usepackage{booktabs}
\usepackage{longtable}
\usepackage{array}
\usepackage{enumitem}
\usepackage{titlesec}
\usepackage{fancyhdr}
\usepackage{setspace}
\usepackage{float}
\usepackage{hyperref}
\usepackage{tikz}
\usepackage{amsmath}
\usepackage{listings}
\usepackage{mdframed}
\usepackage{caption}

% ── Цвета ─────────────────────────────────────────────────
\definecolor{amber}{RGB}{180,110,10}
\definecolor{darkred}{RGB}{140,30,30}
\definecolor{darkgreen}{RGB}{30,100,50}
\definecolor{codebg}{RGB}{245,245,250}
\definecolor{kw}{RGB}{0,80,180}
\definecolor{str}{RGB}{0,120,50}
\definecolor{cmt}{RGB}{120,120,140}
\definecolor{num}{RGB}{160,60,0}

% ── Листинги ──────────────────────────────────────────────
\lstdefinestyle{js}{
  language=JavaScript,
  backgroundcolor=\color{codebg},
  basicstyle=\ttfamily\small,
  keywordstyle=\color{kw}\bfseries,
  stringstyle=\color{str},
  commentstyle=\color{cmt}\itshape,
  numberstyle=\tiny\color{gray},
  numbers=left, numbersep=6pt,
  frame=single, framesep=3pt,
  rulecolor=\color{gray!35},
  breaklines=true, captionpos=b,
  tabsize=2, showstringspaces=false,
  morekeywords={const,let,async,await,import,export,
                from,class,extends,of,true,false},
}

\lstdefinestyle{html}{
  language=HTML,
  backgroundcolor=\color{codebg},
  basicstyle=\ttfamily\small,
  keywordstyle=\color{kw}\bfseries,
  stringstyle=\color{str},
  commentstyle=\color{cmt}\itshape,
  numberstyle=\tiny\color{gray},
  numbers=left, numbersep=6pt,
  frame=single, framesep=3pt,
  rulecolor=\color{gray!35},
  breaklines=true, captionpos=b,
  tabsize=2,
}

\lstdefinestyle{css}{
  language=CSS,
  backgroundcolor=\color{codebg},
  basicstyle=\ttfamily\small,
  keywordstyle=\color{num}\bfseries,
  stringstyle=\color{str},
  commentstyle=\color{cmt}\itshape,
  numberstyle=\tiny\color{gray},
  numbers=left, numbersep=6pt,
  frame=single, framesep=3pt,
  rulecolor=\color{gray!35},
  breaklines=true, captionpos=b,
  tabsize=2,
}

% ── Рамка для важных блоков ───────────────────────────────
\newmdenv[
  backgroundcolor=codebg,
  linecolor=amber,
  linewidth=1pt,
  leftline=true, rightline=false,
  topline=false, bottomline=false,
  innerleftmargin=12pt, innerrightmargin=8pt,
  innertopmargin=6pt, innerbottommargin=6pt,
]{notebox}

% ── Колонтитулы ───────────────────────────────────────────
\pagestyle{fancy}
\fancyhf{}
\fancyhead[L]{\small\textit{Технический отчёт — «Секретный Архив»}}
\fancyhead[R]{\small\thepage}
\fancyfoot[C]{\small\textcolor{gray}{\url{https://runick7.github.io/crypto-site/}}}
\renewcommand{\headrulewidth}{0.4pt}
\renewcommand{\footrulewidth}{0.4pt}

% ── Заголовки ─────────────────────────────────────────────
\titleformat{\section}
  {\large\bfseries\color{darkred}}{\thesection.}{0.7em}{}
  [\vspace{-0.3em}\textcolor{amber}{\rule{\linewidth}{0.6pt}}]
\titleformat{\subsection}{\normalsize\bfseries}{\thesubsection.}{0.7em}{}
\titleformat{\subsubsection}{\normalsize\itshape}{\thesubsubsection.}{0.7em}{}

\setlength{\parindent}{0pt}
\setlength{\parskip}{0.45em}
\onehalfspacing

\hypersetup{
  pdftitle={Технический отчёт — Секретный Архив},
  colorlinks=true, linkcolor=darkred, urlcolor=darkred,
}

% =============================================================
\begin{document}

% ── Титул ────────────────────────────────────────────────────
\begin{titlepage}
  \centering
  \vspace*{1.8cm}
  {\small\textcolor{gray}{ТЕХНИЧЕСКИЙ ОТЧЁТ ПО ВЕБ-ПРОЕКТУ}}\\[1em]

  \begin{tikzpicture}
    \draw[amber, line width=1.5pt] (0,0) rectangle (13,2.2);
    \node at (6.5,1.5) {\LARGE\bfseries\color{darkred} Секретный Архив};
    \node at (6.5,0.65) {\normalsize Интерактивный образовательный сайт по криптографии};
  \end{tikzpicture}

  \vspace{2em}

  \begin{tabular}{rl}
    \textbf{Репозиторий:}  & \url{https://github.com/RuNick7/crypto-site}      \\[0.4em]
    \textbf{Хостинг:}      & \url{https://runick7.github.io/crypto-site/}      \\[0.4em]
    \textbf{Технологии:}   & HTML5 · CSS3 · JavaScript ES2020 (без фреймворков) \\[0.4em]
    \textbf{Объём кода:}   & $\approx$5\,500 строк, 11 исходных файлов         \\[0.4em]
    \textbf{Коммитов:}     & 10                                                \\[0.4em]
    \textbf{Дата:}         & Май 2026                                          \\
  \end{tabular}

  \vfill

  \begin{tabular}{ccc}
    \large\bfseries Оп. VIGENÈRE &
    \large\bfseries Оп. ENIGMA   &
    \large\bfseries Оп. VENONA   \\
    \small 1914–1918 &
    \small 1939–1945 &
    \small 1943–1980 \\
  \end{tabular}

  \vspace{2cm}
\end{titlepage}

\tableofcontents
\newpage

% =============================================================
\section{Цели и общая концепция}

Проект представляет собой статический образовательный сайт, демонстрирующий
три ключевых эпизода истории военной криптографии. Каждая операция реализована
как многошаговый интерактивный урок: пользователь сначала знакомится с принципом
шифра через визуализации и симуляторы, а затем решает фиксированное задание
с проверкой ответа.

\textbf{Ключевые требования к реализации:}
\begin{itemize}[leftmargin=2em]
  \item нет сборщиков, нет фреймворков — только нативный HTML/CSS/JS;
  \item вся криптографическая логика отделена от DOM-кода;
  \item задания имеют однозначно верный ответ с валидацией;
  \item при неверном ответе пользователь получает пошаговое объяснение;
  \item операции можно проходить повторно (кнопка «Начать заново»);
  \item прогресс хранится в \texttt{localStorage}.
\end{itemize}

% =============================================================
\section{Архитектура проекта}

\subsection{Файловая структура}

\begin{verbatim}
crypto-site/
├── index.html            — главная страница (428 строк)
├── vigenere.html         — операция Vigenère (811 строк)
├── enigma.html           — операция Enigma (1072 строки)
├── venona.html           — операция Venona (797 строк)
├── css/
│   ├── base.css          — CSS-переменные, сброс, типографика (258)
│   ├── layout.css        — навигация, сетка, карточки (326)
│   ├── design.css        — компоненты UI (1130)
│   └── animations.css    — keyframe-анимации (253)
└── js/
    ├── ciphers/
    │   ├── vigenere.js   — шифр Виженера (118)
    │   ├── enigma.js     — машина Энигма (197)
    │   └── otp.js        — одноразовый блокнот (75)
    └── ui/
        └── common.js     — бургер-меню, прогресс-бар
\end{verbatim}

\subsection{Принцип разделения логики и представления}

Каждый файл \texttt{js/ciphers/*.js} — чистый ES-модуль без каких-либо
ссылок на DOM. Он экспортирует только функции и классы над строками и числами.
HTML-страницы подключают эти модули через \verb|<script type="module">|
и используют их результаты для построения интерфейса.

\begin{lstlisting}[style=html, caption={Подключение модуля в vigenere.html}]
<script type="module">
import {
  normalize, encrypt, decrypt, encryptSteps,
  kasiskiTrigrams, gcdList, splitColumns,
  guessShift, letterFrequency
} from './js/ciphers/vigenere.js';

// Вся DOM-логика находится здесь
</script>
\end{lstlisting}

\subsection{Дизайн-токены (CSS-переменные)}

Вся визуальная система определена через CSS-переменные в \texttt{base.css},
что обеспечивает единообразие на всех страницах:

\begin{lstlisting}[style=css, caption={Фрагмент дизайн-токенов из base.css}]
:root {
  /* Фон */
  --bg:        #0a0a0f;
  --bg-card:   #0f0f18;
  --bg-input:  #111122;
  --border:    #1e1e2e;

  /* Цвета */
  --amber:     #f5a623;
  --amber-dim: #a06010;
  --green:     #00ff41;
  --green-dim: #00aa2b;
  --red:       #ff3333;
  --white:     #e8e8f0;
  --text-dim:  #666680;

  /* Шрифты */
  --font-title: 'Special Elite', serif;
  --font-mono:  'Share Tech Mono', monospace;
  --btn-font:   'Bebas Neue', sans-serif;

  /* Прочее */
  --radius:     4px;
  --transition: 0.2s ease;
}
\end{lstlisting}

\subsection{Система прогресса}

Прохождение каждой операции сохраняется в \texttt{localStorage}.
Скрипт \texttt{common.js} читает эти значения на каждой странице
и обновляет счётчик «ДОПУСК N/3» и прогресс-бар:

\begin{lstlisting}[style=js, caption={common.js — синхронизация прогресса}]
const ops  = ['vigenere', 'enigma', 'venona'];
const done = ops.filter(o => localStorage.getItem('done_' + o)).length;
if (lvlEl) lvlEl.textContent = done;
if (fill)  fill.style.width  = (done / 3 * 100) + '%';
\end{lstlisting}

% =============================================================
\section{Операция VIGENÈRE: реализация}

\subsection{Математическая основа}

Шифр Виженера — полиалфавитная подстановка с периодическим ключом длины $m$:
\begin{equation}
  c_i = (p_i + k_{i \bmod m}) \bmod 26, \qquad
  p_i = (c_i - k_{i \bmod m} + 26) \bmod 26
\end{equation}

\subsection{Реализация в vigenere.js}

\begin{lstlisting}[style=js, caption={Шифрование и расшифрование (vigenere.js)}]
function encrypt(plaintext, keyword) {
  const pt  = normalize(plaintext);   // убрать не-A..Z
  const key = normalize(keyword);
  if (!key.length) return pt;
  return pt.split('').map((ch, i) => {
    const shift = ALPHA.indexOf(key[i % key.length]);
    return ALPHA[(ALPHA.indexOf(ch) + shift) % 26];
  }).join('');
}

function decrypt(ciphertext, keyword) {
  const ct  = normalize(ciphertext);
  const key = normalize(keyword);
  return ct.split('').map((ch, i) => {
    const shift = ALPHA.indexOf(key[i % key.length]);
    return ALPHA[(ALPHA.indexOf(ch) - shift + 26) % 26];
  }).join('');
}
\end{lstlisting}

Функция \texttt{encryptSteps} возвращает массив объектов
\verb|{plain, key, shift, enc}| для пошаговой визуализации в таблице Виженера.

\subsection{Метод Казиски}

Функция \texttt{kasiskiTrigrams} находит все повторяющиеся триграммы
в шифртексте и вычисляет расстояния между их вхождениями:

\begin{lstlisting}[style=js, caption={Поиск повторяющихся триграмм (vigenere.js)}]
function kasiskiTrigrams(text) {
  const t = normalize(text);
  const seen = {};
  for (let i = 0; i <= t.length - 3; i++) {
    const tri = t.slice(i, i + 3);
    if (!seen[tri]) seen[tri] = [];
    seen[tri].push(i);
  }
  const result = [];
  for (const [tri, positions] of Object.entries(seen)) {
    if (positions.length < 2) continue;
    const distances = [];
    for (let i = 1; i < positions.length; i++)
      distances.push(positions[i] - positions[0]);
    result.push({ tri, positions, distances });
  }
  return result.sort((a, b) => b.positions.length - a.positions.length);
}
\end{lstlisting}

Наибольший общий делитель расстояний вычисляется через \texttt{gcdList}
(итеративный алгоритм Евклида) и выводится как вероятная длина ключа.

\subsection{Частотный анализ}

Функция \texttt{splitColumns} разбивает шифртекст на $m$ колонок, каждая из
которых представляет шифр Цезаря с одним фиксированным сдвигом.
Функция \texttt{guessShift} находит самую частую букву в колонке и предполагает,
что она соответствует \texttt{E} (наиболее частая буква английского языка):

\begin{equation}
  \text{shift} = (\text{pos}(\text{top\_letter}) - \text{pos}(\texttt{E}) + 26) \bmod 26
\end{equation}

\subsection{Задание с фиксированным ответом}

В финальном шаге определён фиксированный перехват:

\begin{lstlisting}[style=js, caption={Определение задания в vigenere.html}]
const CHALLENGE_PLAIN  = 'GERMANADVANCEONPARISALLUNITSMOVEFORWARD';
const CHALLENGE_KEY    = 'WAR';
const CHALLENGE_CIPHER = encrypt(CHALLENGE_PLAIN, CHALLENGE_KEY);
// CHALLENGE_CIPHER вычисляется при загрузке страницы
\end{lstlisting}

\noindent\textbf{Логика валидации:}

\begin{lstlisting}[style=js, caption={Проверка ключа пользователя}]
document.getElementById('revealBtn').addEventListener('click', () => {
  const key = normalize(document.getElementById('revealKey').value);
  if (key === CHALLENGE_KEY) {
    showDecryptedSuccess(key);   // анимация flip + сводка операции
  } else {
    input.classList.add('shake'); // анимация дрожания
    setTimeout(() => input.classList.remove('shake'), 400);
    document.getElementById('wrongKeyMsg').style.display = 'block';
  }
});
\end{lstlisting}

При неверном ответе отображается красный блок с подсказкой и кнопкой
«Узнать ответ с объяснением», которая раскрывает правильный ключ и
пошаговое описание метода взлома.

% =============================================================
\section{Операция ENIGMA: реализация}

\subsection{Класс EnigmaMachine}

Класс реализует полную симуляцию исторической Энигмы I (Wehrmacht/Luftwaffe)
с историческими проводками пяти роторов и двух рефлекторов:

\begin{lstlisting}[style=js, caption={Конструктор EnigmaMachine (enigma.js)}]
class EnigmaMachine {
  constructor(
    rotorNames    = ['I','II','III'],
    reflectorName = 'B',
    startPositions = [0,0,0],
    ringSettings  = [0,0,0],
    plugboardPairs = []
  ) {
    this.rotorDefs = rotorNames.map(n => ROTORS[n]);
    this.reflector = REFLECTORS[reflectorName];
    this.positions = [...startPositions];
    this.rings     = [...ringSettings];
    this.plugboard = buildPlugboard(plugboardPairs);
  }
  // ...
}
\end{lstlisting}

Исторические проводки хранятся как строки из 26 символов — позиция в строке
соответствует входному контакту, символ — выходному:

\begin{lstlisting}[style=js, caption={Исторические данные роторов}]
const ROTORS = {
  I:   { wiring: 'EKMFLGDQVZNTOWYHXUSPAIBRCJ', notch: 'Q' },
  II:  { wiring: 'AJDKSIRUXBLHWTMCQGZNPYFVOE', notch: 'E' },
  III: { wiring: 'BDFHJLCPRTXVZNYEIWGAKMUSQO', notch: 'V' },
  IV:  { wiring: 'ESOVPZJAYQUIRHXLNFTGKDCMWB', notch: 'J' },
  V:   { wiring: 'VZBRGITYUPSDNHLXAWMJQOFECK', notch: 'Z' },
};
\end{lstlisting}

\subsection{Шаг роторов с двойным проскоком}

Реальная Энигма имела конструктивный баг — «двойное проскакивание» (double-stepping
anomaly): средний ротор делал два шага подряд при определённой позиции.
Это поведение воспроизведено точно:

\begin{lstlisting}[style=js, caption={Метод step() с двойным проскоком}]
step() {
  const notches = this.rotorDefs.map(r => r.notch);
  const pos = this.positions;
  const midAtNotch   = ALPHA[pos[1]] === notches[1];
  const rightAtNotch = ALPHA[pos[2]] === notches[2];
  if (midAtNotch) {
    pos[0] = (pos[0] + 1) % 26;   // левый тоже шагает
    pos[1] = (pos[1] + 1) % 26;
  }
  if (rightAtNotch) {
    pos[2] = (pos[2] + 1) % 26;
    if (!midAtNotch) pos[1] = (pos[1] + 1) % 26;
  } else {
    pos[2] = (pos[2] + 1) % 26;
  }
}
\end{lstlisting}

\subsection{Кодирование одной буквы}

Метод \texttt{encodeLetter} возвращает не только результат, но и
полный путь сигнала (\texttt{path}) для визуализации в интерфейсе:

\begin{lstlisting}[style=js, caption={Путь сигнала через машину}]
encodeLetter(letter) {
  const path = [];
  // 1. Штекерная панель (вход)
  const pbIn = this.plugboard[letter] || letter;
  path.push({ stage: 'plugboard_in', in: letter, out: pbIn });

  this.step(); // поворот роторов ДО прохождения сигнала
  let ch = pbIn;

  // 2. Роторы слева направо (прямой ход)
  for (let i = 2; i >= 0; i--) {
    ch = rotorEncode(ch, this.rotorDefs[i].wiring,
                     this.positions[i], this.rings[i]);
    path.push({ stage: `rotor${i+1}_fwd`, in: '...', out: ch });
  }

  // 3. Рефлектор
  ch = this.reflector[ALPHA.indexOf(ch)];
  path.push({ stage: 'reflector', in: '...', out: ch });

  // 4. Роторы справа налево (обратный ход)
  for (let i = 0; i <= 2; i++) {
    const rev = reverseWiring(this.rotorDefs[i].wiring);
    ch = rotorEncode(ch, rev, this.positions[i], this.rings[i]);
    path.push({ stage: `rotor${i+1}_rev`, in: '...', out: ch });
  }

  // 5. Штекерная панель (выход)
  const pbOut = this.plugboard[ch] || ch;
  path.push({ stage: 'plugboard_out', in: ch, out: pbOut });

  return { output: pbOut, path };
}
\end{lstlisting}

\subsection{Криб-драггинг}

Функция \texttt{cribDrag} перебирает все позиции криба в шифртексте
и отбрасывает те, где хотя бы одна буква криба совпадает с буквой шифртекста
(нарушение правила самозашифровки Энигмы):

\begin{lstlisting}[style=js, caption={Алгоритм cribDrag (enigma.js)}]
function cribDrag(ciphertext, crib) {
  const ct = ciphertext.toUpperCase().replace(/[^A-Z]/g,'');
  const cb = crib.toUpperCase().replace(/[^A-Z]/g,'');
  const possible = [], impossible = [];

  for (let offset = 0; offset + cb.length <= ct.length; offset++) {
    let selfEnc = false;
    for (let i = 0; i < cb.length; i++) {
      if (ct[offset + i] === cb[i]) { selfEnc = true; break; }
    }
    selfEnc ? impossible.push(offset) : possible.push(offset);
  }
  return { possible, impossible };
}
\end{lstlisting}

\subsection{Симуляция Бомбы Тьюринга}

Функция \texttt{bombeSearch} перебирает все $26^3 = 17\,576$ стартовых
позиций роторов. Для каждой позиции она создаёт машину, расшифровывает
первые $|\text{crib}|$ символов шифртекста и проверяет, совпадают ли они
с крибом:

\begin{lstlisting}[style=js, caption={Основной цикл bombeSearch (enigma.js)}]
for (let pos = 0; pos < 26*26*26 && tested < maxTests; pos++) {
  const r = Math.floor(pos / 676) % 26;
  const m = Math.floor(pos / 26)  % 26;
  const l = pos % 26;

  const machine = new EnigmaMachine(
    rotorNames, reflectorName, [r, m, l], ringSettings
  );

  let valid = true;
  for (let i = 0; i < crib.length && i < ct.length; i++) {
    const result = machine.encodeLetter(ct[i]);
    // Энигма симметрична: decode(ct) = pt
    if (!result || result.output !== crib[i]) {
      valid = false; break;
    }
  }
  if (valid) results.push({ posStr: ALPHA[r]+ALPHA[m]+ALPHA[l] });
  tested++;
}
\end{lstlisting}

\textbf{Результат для учебного задания:} шифртекст
\texttt{BFEUPHOKXZLSBEFSLEWXHRZATXTLKUAUUMFLWMURBBDXLHM}\\
(немецкая сводка погоды \texttt{WETTERBERICHTNORDSEE...}, роторы III-II-I,
позиция AAA) — бомба находит ровно \textbf{1 совпадение}: позицию \texttt{AAA}.

\begin{table}[H]
  \centering
  \caption{Параметры фиксированного задания Enigma}
  \begin{tabular}{ll}
    \toprule
    Параметр          & Значение \\
    \midrule
    Открытый текст    & \texttt{WETTERBERICHTNORDSEEGEBIETBEWOELKTLEICHTESWINDE} \\
    Ключ (роторы)     & III -- II -- I, рефлектор B, позиция AAA \\
    Шифртекст         & \texttt{BFEUPHOKXZLSBEFSLEWXHRZATXTLKUAUUMFLWMURBBDXLHM} \\
    Криб              & \texttt{WETTER} (первые 6 букв) \\
    Правильный ответ  & Позиция \texttt{AAA} \\
    \bottomrule
  \end{tabular}
\end{table}

% =============================================================
\section{Операция VENONA: реализация}

\subsection{Одноразовый блокнот (otp.js)}

Функции \texttt{otpApply} и \texttt{otpDecrypt} реализуют шифр Вернама
как сложение/вычитание по модулю 26:

\begin{lstlisting}[style=js, caption={OTP шифрование/расшифрование (otp.js)}]
function otpApply(text, key) {
  const t = normalize(text);
  const k = normalize(key);
  return t.split('').map((ch, i) => {
    const ti = ALPHA.indexOf(ch);
    const ki = ALPHA.indexOf(k[i % k.length]);
    return ALPHA[(ti + ki) % 26];
  }).join('');
}
\end{lstlisting}

\subsection{Атака глубины: XOR двух шифртекстов}

Если два сообщения зашифрованы одним ключом $K$, то:
\begin{equation}
  C_1 \oplus C_2 = (P_1 \oplus K) \oplus (P_2 \oplus K) = P_1 \oplus P_2
\end{equation}

Функция \texttt{xorCiphertexts} вычисляет побуквенную разность по модулю 26:

\begin{lstlisting}[style=js, caption={XOR двух шифртекстов (otp.js)}]
function xorCiphertexts(c1, c2) {
  const a = normalize(c1);
  const b = normalize(c2);
  const len = Math.min(a.length, b.length);
  let result = '';
  for (let i = 0; i < len; i++) {
    result += ALPHA[(ALPHA.indexOf(a[i]) - ALPHA.indexOf(b[i]) + 26) % 26];
  }
  return result;
}
\end{lstlisting}

\subsection{Криб-драггинг по XOR}

Имея $P_1 \oplus P_2$ и предполагаемое слово из $P_2$, можно восстановить
соответствующий фрагмент $P_1$:

\begin{lstlisting}[style=js, caption={Криб-драггинг по XOR (otp.js)}]
function cribDragXor(xorResult, crib, offset = 0) {
  const x = normalize(xorResult);
  const c = normalize(crib);
  const recovered = [];
  for (let i = 0; i < c.length && offset + i < x.length; i++) {
    const xi = ALPHA.indexOf(x[offset + i]);
    const ci = ALPHA.indexOf(c[i]);
    recovered.push(ALPHA[(xi + ci) % 26]);
  }
  return recovered.join('');
}
\end{lstlisting}

% =============================================================
\section{Интерфейсные механики}

\subsection{Пошаговая навигация}

Все три операции реализованы как многошаговые мастера. Каждый шаг —
это \verb|<div class="step-panel">|, который показывается или скрывается
через \texttt{display: block/none}. Функция \texttt{goStep(n)}:
\begin{itemize}[leftmargin=2em]
  \item скрывает текущую панель и показывает нужную;
  \item обновляет точки прогресса (классы \texttt{active}, \texttt{done});
  \item плавно прокручивает страницу к началу контента.
\end{itemize}

\subsection{Система анимаций}

Все анимации определены в \texttt{animations.css} как CSS-keyframes.
Применяются через добавление класса из JavaScript:

\begin{table}[H]
  \centering
  \caption{Ключевые анимации проекта}
  \begin{tabular}{lp{8.5cm}}
    \toprule
    \textbf{Keyframe / класс} & \textbf{Применение} \\
    \midrule
    \texttt{flipIn / .flip-char}       & Буквы расшифрованного текста появляются по одной с переворотом по оси X \\
    \texttt{shake / .shake}            & Поле ввода при неверном ключе дрожит 0.4 с \\
    \texttt{fadeUp / .fade-up}         & Панели появляются снизу вверх при переходе между шагами \\
    \texttt{popIn / .pop-in}           & Элементы «выпрыгивают» при появлении \\
    \texttt{glowPulse / .glow-green}   & Пульсирующее зелёное свечение найденного ключа \\
    \texttt{drumRollFast}              & Барабаны Бомбы прыгают при переборе позиций \\
    \texttt{signalBtnPulse}            & Янтарное мерцание кнопки «Проследить сигнал» \\
    \texttt{fall / .rain-char}         & Матричный дождь на главной странице \\
    \bottomrule
  \end{tabular}
\end{table}

\subsection{Стаггер-анимации}

Утилитарный класс \texttt{.stagger} анимирует дочерние элементы
с нарастающей задержкой, создавая эффект последовательного появления:

\begin{lstlisting}[style=css, caption={Утилита .stagger из animations.css}]
.stagger > * { opacity: 0; }
.stagger > *:nth-child(1) { animation: fadeUp 0.5s 0.05s ease forwards; }
.stagger > *:nth-child(2) { animation: fadeUp 0.5s 0.15s ease forwards; }
.stagger > *:nth-child(3) { animation: fadeUp 0.5s 0.25s ease forwards; }
/* ...до 6-го элемента */
\end{lstlisting}

\subsection{Валидация и раскрытие ответа}

Схема обработки ответа пользователя (на примере операции VIGENÈRE):

\begin{center}
\begin{tikzpicture}[
  node distance=1.3cm,
  box/.style={rectangle, rounded corners=3pt, draw=gray!60,
              fill=codebg, minimum width=4.8cm, minimum height=0.75cm,
              align=center, font=\small},
  dec/.style={diamond, draw=amber, fill=amber!10,
              minimum width=2.8cm, minimum height=0.7cm,
              align=center, font=\small},
  arr/.style={->, >=stealth, gray!70, thick}
]
  \node[box] (in)     {Ввод ключа};
  \node[dec, below=0.7cm of in] (chk)    {Ключ = WAR?};
  \node[box, below=0.9cm of chk, xshift=-3.6cm] (err) {.shake\\Сообщение об ошибке};
  \node[box, below=0.9cm of chk, xshift=3.6cm]  (ok)  {flipIn-анимация\\Открытый текст};
  \node[box, below=1.0cm of err] (hint)  {Подсказка:\\частотный анализ};
  \node[box, below=0.9cm of hint] (rev)  {«Узнать ответ»:\\ключ + объяснение};
  \node[box, below=0.9cm of ok]   (done) {localStorage.set\\«Операция завершена»};

  \draw[arr] (in)   -- (chk);
  \draw[arr] (chk)  -- node[left, font=\tiny]{нет} (err);
  \draw[arr] (chk)  -- node[right, font=\tiny]{да} (ok);
  \draw[arr] (err)  -- (hint);
  \draw[arr] (hint) -- (rev);
  \draw[arr] (ok)   -- (done);
\end{tikzpicture}
\end{center}

\subsection{Функция перезапуска операции}

На каждой странице реализована кнопка «Начать заново»,
которая полностью сбрасывает состояние:

\begin{lstlisting}[style=js, caption={Функция restartOperation (vigenere.html)}]
function restartOperation() {
  localStorage.removeItem('done_vigenere');
  // Обновить счётчик допуска в навигации
  const newLvl = ['vigenere','enigma','venona']
    .filter(o => localStorage.getItem('done_' + o)).length;
  document.getElementById('clearanceLevel').textContent = newLvl;
  // Сбросить поля Казиски
  document.getElementById('kasiskiInput').value = DEFAULT_CIPHER;
  document.getElementById('kasiskiResult').style.display = 'none';
  // Сбросить результаты частотного анализа
  document.getElementById('freqColumns').innerHTML = '';
  document.getElementById('freqKeyRecover').style.display = 'none';
  foundKey = '';
  foundKeyLen = 3;
  goStep(0);   // вернуть на первый шаг
}
\end{lstlisting}

% =============================================================
\section{Деплой и хостинг}

\subsection{GitHub Pages}

Проект опубликован через GitHub Pages с автоматическим деплоем
при каждом пуше в ветку \texttt{main}. Статус сборки доступен через API:

\begin{verbatim}
gh api repos/RuNick7/crypto-site/pages --jq '.status'
# "built"
\end{verbatim}

Время пересборки составляет 1–2 минуты. Никаких CI/CD-пайплайнов,
сборщиков или package-менеджеров не требуется: GitHub Pages напрямую
раздаёт статические файлы.

\subsection{ES-модули и CORS}

Из-за политики CORS браузеров файлы с \texttt{import} / \texttt{export}
не работают при открытии через протокол \texttt{file://}. Для локальной
разработки необходим HTTP-сервер:

\begin{verbatim}
python3 -m http.server 8080
# открыть http://localhost:8080
\end{verbatim}

На GitHub Pages проблем нет — файлы раздаются по HTTPS.

% =============================================================
\section{История разработки}

\begin{longtable}{lp{10cm}}
  \toprule
  \textbf{Хэш} & \textbf{Что сделано} \\
  \midrule
  \endhead
  \texttt{5e69bad} & Первичный деплой: все три операции, матричный дождь, хронология \\
  \texttt{33b5feb} & Страницы-досье шифров (\texttt{cipher.html}), исправление хронологии \\
  \texttt{3667ec0} & Консистентность имён, исправление скролла к шагам \\
  \texttt{3523309} & Тултип хронологии: не выходит за границы вьюпорта \\
  \texttt{6ba59a1} & Тултип: исправление позиции при первом наведении \\
  \texttt{484ccd2} & Добавлены итоговые блоки во все финальные шаги \\
  \texttt{25e16ff} & Фиксированные задания, валидация ключей, объяснение ответа, кнопки перезапуска \\
  \texttt{fbdbea1} & Немецкое сообщение WETTER, панель перевода, кнопка «Проследить сигнал» \\
  \texttt{c499ff5} & Исправление SVG: текст «Именно это использовал Тьюринг!» выходил за рамку \\
  \texttt{c52bd24} & Исправление \texttt{bombeSearch}: возвращал 2000 совпадений вместо одного; убрана дублирующая метка \\
  \bottomrule
\end{longtable}

% =============================================================
\section{Заключение}

Проект реализует статический образовательный сайт без каких-либо внешних
зависимостей. Вся криптографическая логика написана с нуля: шифр Виженера
с методами Казиски и частотного анализа, полная симуляция машины Энигма I
с историческими проводками и двойным проскоком роторов, атака глубины
на OTP и криб-драггинг по XOR.

Наиболее нетривиальной задачей оказалась симуляция Бомбы Тьюринга:
первоначальная реализация проверяла самозашифровку шифртекста
(почти всегда проходила → 2\,000 ложных совпадений). После исправления
алгоритм правильно сравнивает расшифрованный текст с крибом и возвращает
ровно одно верное совпадение.

\vspace{1.5em}
\begin{center}
  \textcolor{gray}{\rule{10cm}{0.4pt}}\\[0.4em]
  \small\textcolor{gray}{
    \url{https://runick7.github.io/crypto-site/} \quad·\quad
    \url{https://github.com/RuNick7/crypto-site}
  }
\end{center}

\end{document}
